\poczatekwykladu{9}

\section{Niezależność całki od zapisu}
\label{sec:niezaleznosc-calki-od-zapisu}

W dalszej części rozpatrujemy funkcje nieujemne. Przez
$\mathrm{PMN}$ oznaczamy rodzinę prostych mierzalnych funkcji
nieujemnych na $X$. Ich całki są określone i należą do $[0,+\infty]$.
Nadal stosujemy konwencję $0\cdot(+\infty)=0$.

\begin{twierdzenie}[Niezależność od przedstawienia]
\label{tw:calka-niezalezna-od-przedstawienia}
Niech
\[
  f=\sum_{i=1}^{m}a_i\mathbf{1}_{A_i}
   =\sum_{j=1}^{n}b_j\mathbf{1}_{B_j},
\]
gdzie $A_i,B_j\in\mathcal{F}$ oraz $a_i,b_j\geq0$.
Zbiory w tych przedstawieniach nie muszą być rozłączne.
Wówczas
\begin{equation}
  \sum_{i=1}^{m}a_i\mu(A_i)
  =\sum_{j=1}^{n}b_j\mu(B_j)
  =\int_X f\,d\mu.
  \label{eq:calka-dowolny-zapis-prostej}
\end{equation}
\end{twierdzenie}

\begin{proof}
Najpierw rozbijemy rodziny zbiorów na rozłączne części.
Oznaczmy $A_i^1=A_i$ i $A_i^0=A_i^c$. Dla
$k=(k_1,\ldots,k_m)\in\{0,1\}^m$ połóżmy
\[
  A_k=\bigcap_{i=1}^{m}A_i^{k_i}.
\]
Zbiory $A_k$ są mierzalne, parami rozłączne i pokrywają $X$.
Istotnie, każdy punkt ma jednoznacznie określony ciąg przynależności
do zbiorów $A_1,\ldots,A_m$. Niektóre części $A_k$ mogą być puste.
Analogicznie, dla $\ell=(\ell_1,\ldots,\ell_n)\in\{0,1\}^n$
definiujemy
\[
  B_\ell=\bigcap_{j=1}^{n}B_j^{\ell_j},
  \qquad B_j^1=B_j,\quad B_j^0=B_j^c.
\]
Rodzina zbiorów $A_k\cap B_\ell$ jest wspólnym rozłącznym podziałem
$X$ (rysunek~\ref{fig:wspolny-podzial-prostej}).

\begin{figure}[htbp]
  \centering
  \includegraphics[width=0.96\textwidth]{wspolny-podzial-prostej.pdf}
  \caption{Rozbicie rodzin zbiorów i ich wspólny podział: części o ustalonej przynależności do obu rodzin.}
  \label{fig:wspolny-podzial-prostej}
\end{figure}

Na zbiorze $A_k$ pierwszy zapis funkcji ma stałą wartość
\[
  \alpha_k=\sum_{i=1}^{m}k_i a_i,
\]
a na $B_\ell$ drugi zapis ma stałą wartość
$\beta_\ell=\sum_{j=1}^{n}\ell_j b_j$.
Jeżeli $A_k\cap B_\ell\neq\varnothing$, to
$\alpha_k=\beta_\ell$, ponieważ oba zapisy przedstawiają tę samą funkcję.

Dla każdego $i$ mamy rozłączny rozkład
$A_i=\bigcup_{k:\,k_i=1}A_k$. Ze skończonej addytywności miary
i przestawienia skończonych sum nieujemnych składników otrzymujemy
\begin{align*}
  \sum_{i=1}^{m}a_i\mu(A_i)
  &=\sum_{i=1}^{m}\sum_{k:\,k_i=1}a_i\mu(A_k)\\
  &=\sum_{k\in\{0,1\}^m}
      \left(\sum_{i=1}^{m}k_i a_i\right)\mu(A_k)
   =\sum_k\alpha_k\mu(A_k).
\end{align*}
Następnie rozbijamy każdy $A_k$ na części $A_k\cap B_\ell$:
\begin{align*}
  \sum_k\alpha_k\mu(A_k)
  &=\sum_{k,\ell}\alpha_k\mu(A_k\cap B_\ell)\\
  &=\sum_{k,\ell}\beta_\ell\mu(A_k\cap B_\ell)\\
  &=\sum_\ell\beta_\ell\mu(B_\ell)
   =\sum_{j=1}^{n}b_j\mu(B_j).
\end{align*}
W drugiej równości korzystamy z $\alpha_k=\beta_\ell$ na niepustych
częściach podziału; dla części pustych oba składniki są zerowe.
Ostatnia równość jest analogiczna do obliczenia dla rodziny $A_i$.
Porównując dowolny zapis z postacią kanoniczną, otrzymujemy również
równość z całką z definicji~\ref{def:calka-prostej}.
\end{proof}

\section{Własności całki funkcji prostej}
\label{sec:wlasnosci-calki-prostej}

\begin{twierdzenie}
\label{tw:wlasnosci-calki-prostej}
Dla $f,g\in\mathrm{PMN}$ i $c\geq0$ zachodzą następujące własności:
\begin{enumerate}
  \item $\displaystyle\int_X f\,d\mu\geq0$;
  \item $\displaystyle\int_X cf\,d\mu=c\int_X f\,d\mu$;
  \item $\displaystyle\int_X(f+g)\,d\mu
          =\int_X f\,d\mu+\int_X g\,d\mu$;
  \item jeśli $f\leq g$, to
        $\displaystyle\int_X f\,d\mu\leq\int_X g\,d\mu$.
\end{enumerate}
\end{twierdzenie}

\begin{proof}
\emph{1)} Niech $f=\sum_{i=1}^{m}a_i\mathbf{1}_{A_i}$ będzie postacią
kanoniczną. Z $f\geq0$ wynika $a_i\geq0$, a więc wszystkie składniki
$a_i\mu(A_i)$ są nieujemne.

\emph{2)} Mamy $cf=\sum_{i=1}^{m}ca_i\mathbf{1}_{A_i}$.
Z twierdzenia~\ref{tw:calka-niezalezna-od-przedstawienia} wynika
\[
  \int_X cf\,d\mu
  =\sum_{i=1}^{m}ca_i\mu(A_i)
  =c\sum_{i=1}^{m}a_i\mu(A_i)
  =c\int_X f\,d\mu.
\]
Dla $c=0$ obie strony są zerowe, także gdy całka z $f$ jest nieskończona.

\emph{3)} Niech $f=\sum_{i=1}^{m}a_i\mathbf{1}_{A_i}$ i
$g=\sum_{j=1}^{n}b_j\mathbf{1}_{B_j}$ będą postaciami kanonicznymi.
Zbiory $A_i\cap B_j$ tworzą rozłączny podział $X$, przy czym
\[
  f+g=\sum_{i=1}^{m}\sum_{j=1}^{n}
       (a_i+b_j)\mathbf{1}_{A_i\cap B_j}.
\]
Ten zapis nie musi być kanoniczny, lecz możemy obliczyć z niego całkę:
\begin{align*}
  \int_X(f+g)\,d\mu
  &=\sum_{i=1}^{m}\sum_{j=1}^{n}(a_i+b_j)\mu(A_i\cap B_j)\\
  &=\sum_{i=1}^{m}a_i\sum_{j=1}^{n}\mu(A_i\cap B_j)
    +\sum_{j=1}^{n}b_j\sum_{i=1}^{m}\mu(A_i\cap B_j)\\
  &=\sum_{i=1}^{m}a_i\mu(A_i)+\sum_{j=1}^{n}b_j\mu(B_j)\\
  &=\int_X f\,d\mu+\int_X g\,d\mu.
\end{align*}
Użyliśmy rozłączności poziomic oraz równości
$\bigcup_j B_j=X$ i $\bigcup_i A_i=X$.

\emph{4)} Jeżeli $f\leq g$, to $g-f\in\mathrm{PMN}$. Z punktów 1) i 3)
otrzymujemy
\[
  \int_X g\,d\mu
  =\int_X f\,d\mu+\int_X(g-f)\,d\mu
  \geq\int_X f\,d\mu.
\]
Nie odejmujemy tu całek, które mogłyby być nieskończone.
\end{proof}

\section{Całka po zbiorze mierzalnym}
\label{sec:calka-po-zbiorze}

\begin{definicja}
\label{def:calka-prostej-po-zbiorze}
Dla $E\in\mathcal{F}$ i $f\in\mathrm{PMN}$ definiujemy
\begin{equation}
  \int_E f\,d\mu=\int_X f\mathbf{1}_E\,d\mu.
  \label{eq:calka-po-zbiorze}
\end{equation}
Przyjmujemy konwencję $\int f\,d\mu=\int_X f\,d\mu$.
\end{definicja}

\begin{fakt}
\label{fakt:wlasnosci-calki-po-zbiorze}
Własności z twierdzenia~\ref{tw:wlasnosci-calki-prostej} zachodzą także
dla całki po zbiorze $E\in\mathcal{F}$. W szczególności
\[
  \int_E(f+g)\,d\mu=\int_E f\,d\mu+\int_E g\,d\mu.
\]
\end{fakt}

\begin{proof}
Stosujemy odpowiednie własności do funkcji $f\mathbf{1}_E$
i $g\mathbf{1}_E$. Na przykład
$(f+g)\mathbf{1}_E=f\mathbf{1}_E+g\mathbf{1}_E$.
\end{proof}

\section{Całka funkcji nieujemnej: etap drugi}
\label{sec:calka-nieujemnej-etap-drugi}

Przez $\mathrm{MN}$ oznaczamy rodzinę mierzalnych funkcji nieujemnych
$f\colon X\to[0,+\infty]$.

\begin{definicja}[Całka funkcji nieujemnej]
\label{def:calka-nieujemnej}
Dla $f\in\mathrm{MN}$ definiujemy
\begin{equation}
  \int_X f\,d\mu
  =\sup\left\{\int_X s\,d\mu:
       s\in\mathrm{PMN},\ s\leq f\right\}.
  \label{eq:calka-nieujemnej-supremum}
\end{equation}
Nierówność $s\leq f$ ma zachodzić w każdym punkcie $X$.
\end{definicja}

Całkę wyznaczamy zatem za pomocą całek funkcji prostych leżących pod $f$.
Rodzina w definicji jest niepusta, ponieważ zawiera funkcję zerową.
Supremum może być równe $+\infty$.
Schematy przedstawia rysunek~\ref{fig:calka-nieujemnej}.

\begin{figure}[htbp]
  \centering
  \includegraphics[width=0.96\textwidth]{calka-nieujemnej.pdf}
  \caption{Po lewej: funkcje proste nieujemne leżące pod $f$. Po prawej: ograniczenie całkowania do zbioru $E$.}
  \label{fig:calka-nieujemnej}
\end{figure}

\begin{uwaga}[Zgodność z etapem pierwszym]
\label{uw:zgodnosc-etapow-calki}
Jeżeli $f\in\mathrm{PMN}$, definicja~\ref{def:calka-nieujemnej}
daje tę samą wartość co definicja~\ref{def:calka-prostej}.
Istotnie, dla $s\in\mathrm{PMN}$ spełniającej $s\leq f$ monotoniczność
całki funkcji prostej daje $\int_X s\,d\mu\leq\int_X f\,d\mu$.
Z drugiej strony w supremum możemy wybrać $s=f$.
\end{uwaga}

\begin{twierdzenie}[Własności całki funkcji nieujemnej]
\label{tw:wlasnosci-calki-nieujemnej}
Dla $f,g\in\mathrm{MN}$ i $c\geq0$ mamy:
\begin{enumerate}
  \item $\displaystyle\int_X f\,d\mu\geq0$;
  \item $\displaystyle\int_X cf\,d\mu=c\int_X f\,d\mu$;
  \item $\displaystyle\int_X(f+g)\,d\mu
          =\int_X f\,d\mu+\int_X g\,d\mu$;
  \item jeśli $f\leq g$, to
        $\displaystyle\int_X f\,d\mu\leq\int_X g\,d\mu$.
\end{enumerate}
\end{twierdzenie}

\begin{proof}[Uzasadnienie punktów 1), 2) i 4)]
\emph{1)} Wszystkie całki występujące pod supremum w
\eqref{eq:calka-nieujemnej-supremum} są nieujemne.

\emph{2)} Dla $c>0$ warunek $s\in\mathrm{PMN}$ i $s\leq cf$ jest
równoważny warunkowi $s=c\widetilde{s}$, gdzie
$\widetilde{s}\in\mathrm{PMN}$ i $\widetilde{s}\leq f$.
Korzystając z jednorodności całki funkcji prostej, dostajemy
\begin{align*}
  \int_X cf\,d\mu
  &=\sup\left\{\int_X s\,d\mu:s\in\mathrm{PMN},\ s\leq cf\right\}\\
  &=\sup\left\{c\int_X\widetilde{s}\,d\mu:
             \widetilde{s}\in\mathrm{PMN},\ \widetilde{s}\leq f\right\}\\
  &=c\sup\left\{\int_X\widetilde{s}\,d\mu:
             \widetilde{s}\in\mathrm{PMN},\ \widetilde{s}\leq f\right\}
   =c\int_X f\,d\mu.
\end{align*}
Dla $c=0$ funkcja $cf$ jest zerowa i obie strony są zerowe,
przy przyjętej konwencji $0\cdot(+\infty)=0$.

\emph{4)} Jeśli $f\leq g$, każda funkcja prosta nieujemna $s\leq f$
spełnia także $s\leq g$. Supremum dla $f$ bierzemy więc po podrodzinie
rodziny użytej dla $g$, co daje żądaną nierówność.
\end{proof}
